\chapter[Ein Geleitwort zur Ausfahrt]{Ein Geleitwort zur Ausfahrt}
\quellseite{11}
Kein Unglück kann sich ereignen, nichts so Trauriges oder Bitteres über uns kommen, das nicht, sobald die heilige Schrift uns aufnimmt, entweder schnell verscheucht oder leicht erduldet würde.

Wie oft, theuerster Jesu, wird es bei mir Abend! Wie oft folgt auf das Tageslicht einer kleinen Tröstung ein unerträglicher Schmerz gleich nächtlichem Dunkel! Es wird mir alles zuwider. Zur Last wird mir alles, was ich sehe. Es spricht jemand und kaum höre ich es. Es klopft jemand an und kaum merke ich es. Hart wie ein Stein wird mein Herz. Die Zunge stockt und die Thränenbäche der Augen vertrocknen. Wie dann? Ich gehe wenigstens auf Betrachtungen aus, ich suche in der heiligen Schrift. Und siehe, unvermuthet, plötzlich kommt mir deine Gnade entgegen, theurer Jesu! Mit ihrem Lichte verscheucht sie das Dunkel, vertreibt sie den Unmuth, bricht sie die Hülle. Bald folgen auf Seufzer Thränen, und in die Thränen mischt sich die himmlische Freude.

O die Armen, die, wenn sie Traurigkeit drückt, nicht diese Auen betreten, um Freude zu gewinnen! (Alexis, † 1166.)
